% Course study notes. Copy to runs/<course>/notes/notes.tex and replace every
% <<...>> marker; build_notes.py refuses to pass while any marker remains.
% Each outline unit is its own file under notes/units/ (one \section each).
\documentclass[11pt]{article}
\usepackage[letterpaper,margin=1in]{geometry}
\usepackage{amsmath,amssymb,amsthm,mathtools,bm}
\usepackage{booktabs,tabularx,array,longtable}
\usepackage[shortlabels]{enumitem}
\usepackage{tikz}
\usetikzlibrary{arrows.meta,positioning,calc,decorations.pathreplacing,patterns,shapes.geometric,matrix}
\usepackage{pgfplots}
\pgfplotsset{compat=1.18}
\usepackage{xcolor}
\usepackage{microtype}
\usepackage{fancyhdr}
\usepackage[hidelinks]{hyperref}

\setlength{\parskip}{0.45em}
\setlength{\parindent}{0pt}
\setlist{itemsep=0.15em,topsep=0.3em}
\pagestyle{fancy}
\fancyhf{}
\fancyhead[L]{\small\nouppercase{\leftmark}}
\fancyhead[R]{\small <<SHORT COURSE NAME>>}
\fancyfoot[C]{\small\thepage}
\renewcommand{\headrulewidth}{0.3pt}

% Numbered environments share one counter per section: Definition 2.1, Theorem 2.2, ...
\theoremstyle{definition}
\newtheorem{definition}{Definition}[section]
\newtheorem{example}[definition]{Example}
\newtheorem{exercise}[definition]{Exercise}
\theoremstyle{plain}
\newtheorem{theorem}[definition]{Theorem}
\newtheorem{proposition}[definition]{Proposition}
\newtheorem{lemma}[definition]{Lemma}
\newtheorem{corollary}[definition]{Corollary}
\theoremstyle{remark}
\newtheorem{remark}[definition]{Remark}
\newtheorem*{trap}{Common trap}
\newenvironment{solution}{\begin{proof}[Solution]}{\end{proof}}

% Course-specific notation (add what the course uses; keep it consistent).
\newcommand{\R}{\mathbb{R}}
\newcommand{\N}{\mathbb{N}}
\newcommand{\Z}{\mathbb{Z}}
\newcommand{\Q}{\mathbb{Q}}
\newcommand{\C}{\mathbb{C}}

\begin{document}

\begin{titlepage}
  \raggedright
  \vspace*{2.2cm}
  {\Huge\bfseries <<COURSE TITLE>>\par}
  \vspace{0.8em}
  {\Large <<COURSE CODE AND INSTITUTION>>\par}
  \vspace{0.4em}
  {\large Study notes\par}
  \vfill
  \begin{flushright}
    \begin{minipage}{0.46\textwidth}
      \raggedleft\hyphenpenalty=10000\exhyphenpenalty=10000
      {\itshape ``<<A SHORT, VERIFIED QUOTE BY A PIONEER OF THIS FIELD>>''\par}
      \vspace{0.5em}
      --- <<NAME>>\par
    \end{minipage}
  \end{flushright}
  \vfill
  {\small <<MONTH YEAR>>\par}
\end{titlepage}

\tableofcontents
\clearpage

\section*{About these notes}
<<TWO OR THREE SENTENCES: WHAT THE COURSE COVERS, WHAT A READER SHOULD KNOW FIRST, AND HOW THE SECTIONS BUILD ON EACH OTHER.>>

\subsection*{Notation}
\begin{longtable}{@{}p{0.28\textwidth}p{0.66\textwidth}@{}}
\toprule
Symbol & Meaning \\
\midrule
<<SYMBOL>> & <<MEANING>> \\
\bottomrule
\end{longtable}
\clearpage

% One file per outline unit, in course order.
\input{units/u01}

\appendix
\section{Formula sheet}
<<THE KEY RESULTS OF THE COURSE IN ONE PLACE, GROUPED BY SECTION, EACH WITH A POINTER TO WHERE IT IS DERIVED.>>

\end{document}
